\documentclass[11pt]{article}
\usepackage[margin=1in]{geometry}
\usepackage{amsmath,amssymb,amsthm,mathtools}
\usepackage[hidelinks]{hyperref}
\allowdisplaybreaks[2]
\setlength{\emergencystretch}{3em}
\newtheorem{theorem}{Theorem}[section]
\newtheorem{lemma}[theorem]{Lemma}
\newtheorem{proposition}[theorem]{Proposition}
\theoremstyle{remark}\newtheorem{remark}[theorem]{Remark}
\newcommand{\dd}{\,d}
\newcommand{\E}{\mathbb E}
\newcommand{\Rhalf}{\mathcal R}
\newcommand{\ones}{\boldsymbol 1}
\title{Strong $M$-clock passage: a fixed-cutoff obstruction}
\author{Analytical increment for review}
\date{Version 1; October 6, 2026}
\begin{document}\maketitle
\begin{abstract}
The proposed uniform $e^{o(1)}$ transfer by the linear centered
$M$-clock multiplier does not follow from the overlap data.
At fixed overlap offset $C_m$, its outer Q1 labels have scaled
displacement comparable to $e^{-dC_m}$, not $o(1)$.
We prove, for the initialized original flow on $[t_m,t_m+1]$, a
strict positive correction of that order to the logarithm of the
linear shape multiplier. The corresponding normalized-tangent
coefficient has an integral bounded away from zero as $q\to0$.
An exact cooperative leading-field example retains this correction
all the way to $M=1/2$. This example is a test of the proposed
leading comparison, not a substitute initialized population.
We give the corrected handoff normalization and stop at the
central tangent obstruction. No original-flow half-mass theorem
is asserted.
\end{abstract}

\section{Scope and the overlap data}
All new labels have prefix \texttt{an02smcpv1:}.
Use the initialized model and fixed right-half-circle ancestry
$\Rhalf=[-\pi/2,\pi/2]$ of OWPT, with label variable $\beta$,
initial radial mass $q^{10}$, and $\rho\in[13/20,7/10]$.
Write
\[
 \lambda=\rho^2,\quad \omega=(1-\lambda)/2,\quad
 a=\rho K(\rho a),\quad \alpha=\lambda\Phi(\rho a)/2,
 \quad d=1/2-\alpha,\quad s=d/\omega,
\]
where $K(z)=\phi(z)+z\Phi(z)$.
The source bound $\alpha<7497/50000<3/20$ implies
\begin{equation}\label{an02smcpv1:rate}
 r_0(M):=-\frac{1-M}{2}+\alpha=-d+M/2<-1/10
 \qquad (0\le M\le1/2).
\end{equation}
Fix OWPT's sufficiently large, \emph{$q$-independent} $C_m$ and put
\begin{align}
 t_m&=(2/\omega)\log(1/q)+C_m,&
 \sigma&=10-2/\omega,& \zeta&=\sigma-1,\nonumber\\
 E_m&=e^{-dC_m},& \epsilon_m&=E_mq^s,&
 M(t)&=\int_{\Rhalf}m(t,\beta)\frac{\dd\beta}{2\pi}.
 \label{an02smcpv1:scales}
\end{align}
Here $110/51\le\sigma\le710/231$ and $\zeta\ge59/51$.
The same-field reference is always the actual label
$\beta_c=qa_q$, where $a_q=a+O(q^2)$ is PROFILE's exact
target-only root. Set
\[
 h_x(t,\beta)=\frac{\theta(t,\beta)-\theta(t,\beta_c)}q,
 \qquad
 h_y(t,\beta)=\frac{\psi(t,\beta)-\psi(t,\beta_c)}q,
 \qquad
 \widehat h(t,\beta)=\widehat\theta(t,\beta)/q-a_q.
\]
Hats denote the exact target-only flow.
OWPT compares $h_x,h_y$ to the \emph{nonlinear} exact
$\widehat h$, not to its linearization at $a_q$.

Apply that already proved result with the fixed terminal offset
$C_m+1$. Its estimates then hold throughout $[t_m,t_m+1]$:
\begin{align}
 h_x(t,\beta)&=\widehat h(t,\beta)e^{O(q^\zeta)},&
 h_y(t,\beta)&=\widehat h(t,\beta)e^{O(q^\zeta)},
 \label{an02smcpv1:shorttransfer}\\
 m(t,\beta)&=\widehat m(t,\beta)e^{O(q^\sigma)},&
 S(t)&\le q^{10}e^{(1+2q^2)t}=O(q^\sigma).
 \label{an02smcpv1:shortmass}
\end{align}
All constants are uniform in the label and in $\rho$ for fixed
$C_m$. This fixed-length extension uses exactly the existing
overlap proof; it is not a continuation into $S\gtrsim q$.
The actual reference is not reset to $a_q$.

\section{A closed original-flow obstruction on the first unit interval}
Choose PROFILE's common chart constant $L$, enlarged once if
needed. Its proved expansion is
\begin{equation}\label{an02smcpv1:scalar}
 f_q(x):=v_q(qx)/q=f_0(x)+O_{C^2}(q^2),\qquad
 f_0(x)=\tfrac12[\rho K(\rho x)-x],\qquad
 f_0''(x)=\tfrac12\rho^3\phi(\rho x)>0.
\end{equation}
Define the analytically specified positive constant
\begin{equation}\label{an02smcpv1:curvature}
 b=\frac14\min_{\substack{13/20\le\rho\le7/10\\|x|\le L+2}}
                \rho^3\phi(\rho x)>0.
\end{equation}
Thus $f_q''\ge b$ on the relevant positive chart, for small $q$.
Also $-1\le f_q'\le-\omega/2$ there, and
$f_q'(a_q)=-d+O(q^2)$, by the same source expansion.

Let $c_P>0$ be PROFILE's uniform lower comparison constant in
its Q1 displacement law. For the positive-measure initial band
\[
 I_q=\{\beta\in[\pi/4,\pi/2]:0\le\cos\beta\le q\},
\]
the overlap law gives
\begin{equation}\label{an02smcpv1:outerlower}
 \widehat h(t_m,\beta)\ge c_*E_m,\qquad c_*:=c_P/4>0.
\end{equation}
Indeed $(\cos\beta+q)/q\le2$ and
$s<100/51<2$. In particular this lower bound is not an
assertion only about the single endpoint label.

For positive masses define the proposed linear multiplier
\begin{equation}\label{an02smcpv1:multiplier}
 \mathcal G(M_1,M_0)=
 (M_1/M_0)^{-d}\left(\frac{1-M_0}{1-M_1}\right)^\alpha.
\end{equation}

\begin{theorem}[Nonvanishing short-interval correction; original flow]
\label{an02smcpv1:shortobstruction}
Put $M_0=M(t_m)$ and $M_+=M(t_m+1)$. For sufficiently small $q$,
\begin{equation}\label{an02smcpv1:shortclock}
 \log(M_+/M_0)=1+O(q^2+q^\sigma).
\end{equation}
There is $k>0$, independent of $q$, $C_m$, and LCRC's guards,
such that, for every $\beta\in I_q$ and $j=x,y$,
\begin{equation}\label{an02smcpv1:strictshape}
 \log\frac{h_j(t_m+1,\beta)}
 {h_j(t_m,\beta)\mathcal G(M_+,M_0)}
 \ge k E_m>0.
\end{equation}
The small-$q$ threshold may depend on the fixed $C_m$.

Let $\mathcal A_q(t,\beta)$ be the receiver Jacobian of the actual field
in the scaled coordinates $(\theta/q,\psi/q)$, evaluated along
the label $\beta$. Set
\[
 J_{\rm sh}(M)=
 \begin{pmatrix}-(1-M)/2&\alpha\\\alpha&-(1-M)/2\end{pmatrix}.
\]
Possibly decreasing $k$, one has
\begin{equation}\label{an02smcpv1:failedintegral}
 \int_{t_m}^{t_m+1}
 \sup_{\beta\in\Rhalf}\|\mathcal A_q(t,\beta)-J_{\rm sh}(M(t))\|_\infty\dd t
 \ge k E_m.
\end{equation}
The matrix norm is the maximum absolute row sum.
\end{theorem}
\begin{proof}
For $\beta\in I_q$, write $h(u)=\widehat h(t_m+u,\beta)$.
It stays positive by scalar uniqueness and obeys
\[
 h'=f_q(a_q+h),\qquad h(0)e^{-u}\le h(u)\le h(0)
 \quad(0\le u\le1).
\]
The derivative bounds above prove these inequalities, within
PROFILE's invariant chart. Taylor's formula with integral remainder
therefore gives
\begin{align}
 \log\frac{h(1)}{h(0)}+d
 &=\int_0^1\left(\frac{f_q(a_q+h(u))}{h(u)}+d\right)\dd u
 \nonumber\\
 &\ge\frac b2\int_0^1h(u)\dd u-Cq^2
 \ge\frac{bc_*}{2}(1-e^{-1})E_m-Cq^2.
 \label{an02smcpv1:convexdefect}
\end{align}
Division is only by strictly positive displacements. Equation
\eqref{an02smcpv1:shorttransfer} changes the left side by
$O(q^\zeta)$ when $h$ is replaced by either original displacement.

For completeness, the mass statement needs no future logistic
contract. On the target-only chart $|\widehat\theta|\le qL$,
\[
 (\log\widehat m)'=2g(\widehat\theta),\qquad
 2g(\vartheta)=
 \E[(\cos\vartheta+qZ)_+^2]
 +\E[(\rho\sin\vartheta+qZ)_+^2]=1+O_L(q^2).
\]
For the first expectation subtract its negative-part second
moment from $\cos^2\vartheta+q^2$; that moment is at most $q^2$
when $\cos\vartheta\ge0$. The second expectation is $O_L(q^2)$.
Integrate over the unit interval, then use
\eqref{an02smcpv1:shortmass} at its two endpoints. The resulting
pointwise bounds integrate over the fixed positive label measure,
proving \eqref{an02smcpv1:shortclock}. Since $M_0,M_+=O(q^\sigma)$,
\[
 \log\mathcal G(M_+,M_0)=-d+O(q^2+q^\sigma).
\]
This and \eqref{an02smcpv1:convexdefect} prove
\eqref{an02smcpv1:strictshape}, for example with
$k=bc_*(1-e^{-1})/4$ after decreasing $q$.

For the tangent assertion, OWPT's receiver-Jacobian lemma gives
\[
 \mathcal A_q(t,\beta)=
 \begin{pmatrix}-g&\ell\\\ell&-g\end{pmatrix}
 +O_{\ell^\infty}(q^\zeta),\qquad
 \ell-g=f_q'(a_q+\widehat h(t,\beta)),\quad \ell\ge0.
\]
Scaling both coordinates and their velocities by the same $q$
does not change this Jacobian. In particular each component of
$(\mathcal A_q-J_{\rm sh})\ones$ is at least
\[
 b\widehat h(t,\beta)-C(q^2+q^\zeta)-M(t)/2.
\]
Use $M(t)=O(q^\sigma)$, the preceding lower bound for $h(u)$,
and $\|E\|_\infty\ge\|E\ones\|_\infty$.
Integration for any $\beta\in I_q$ proves
\eqref{an02smcpv1:failedintegral}. All errors vanish with $q$;
$E_m$ does not. No finite-noise cooperativity assumption was used.
\end{proof}

\section{Why cooperativity cannot remove this correction}
The failing step in the proposed tangent argument is
\begin{equation}\label{an02smcpv1:requiredfalse}
 \int_{t_m}^{t_{1/2}}\sup_{\beta\in\Rhalf}
 \|\mathcal A_q(t,\beta)-J_{\rm sh}(M(t))\|_\infty\dd t=o(1),
\end{equation}
or an $o(1)$ integral of its row-sum errors that would give the
same uniform normalized-tangent comparison. If the half-mass
hit exists, it is later than $t_m+1$ by
\eqref{an02smcpv1:shortmass}; hence
\eqref{an02smcpv1:failedintegral} contradicts
\eqref{an02smcpv1:requiredfalse}. Even the row sums already
have a positive error on the first unit interval. Allowing a
different Metzler off-diagonal splitting cannot change those
row sums after normalization by the specified scalar multiplier.

This is not merely a poor use of an absolute-value estimate.
The following exact leading calculation shows the nonlinear
effect can persist to the requested endpoint under perfect
cooperativity, zero outside mass, and no finite-$q$ error.

\begin{proposition}[Exact leading-field test at half mass]
\label{an02smcpv1:leadingtest}
In P's resident leading field $\mu_s=M\delta_{(a,a)}$, $N=0$,
let a same-field test characteristic start at
$(a+h_0,a+h_0)$, where $0<h_0\le1/4$ is fixed, and let the
reference be $(a,a)$. With $M'=M(1-M)$ and $M(0)=M_0\to0$,
write $T$ for the first hit of $M=1/2$.
The test characteristic remains aligned and in the fixed chart
$a\le x=y\le a+1/4$. For constants $k_->0$ and $k_+=1/8$,
\begin{equation}\label{an02smcpv1:leadingratio}
 2k_-h_0(1-e^{-1/2})
 \le\log\frac{x(T)-a}{h_0\mathcal G(1/2,M_0)}
 \le20k_+h_0
 \qquad(T\ge1).
\end{equation}
Thus the ratio is not $e^{o(1)}$ as $M_0\to0$.
\end{proposition}
\begin{proof}
P's overlap kernel has
$F(\rho x,\rho a)=K(\rho a)$ for $x\ge a$.
Insert this and $Y=Ma$ into its strong receiver field. At
$x=y=a+h$, the two velocities agree and
\begin{equation}\label{an02smcpv1:nonlinearleading}
 h'=r_0(M)h+R(h),\qquad
 R(h)=\frac\rho2\{K(\rho(a+h))-K(\rho a)
                         -\rho\Phi(\rho a)h\}.
\end{equation}
Strict convexity and Taylor's formula give
$k_-h^2\le R(h)\le k_+h^2$ on $0\le h\le1/4$, where
\[
 k_-:=\frac14\min_{\substack{13/20\le\rho\le7/10\\0\le v\le1/4}}
              \rho^3\phi(\rho(a_\rho+v))>0.
\]
The upper choice $k_+=1/8$ uses only $\rho<1$ and
$\phi(0)<1/2$. By \eqref{an02smcpv1:rate},
$h'\le(-1/10+1/32)h\le-h/20$, while $h'\ge-h/2$.
These inward inequalities prove the chart and
$h_0e^{-t/2}\le h(t)\le h_0e^{-t/20}$ up to $T$.
The logistic substitution, already used in P, gives
$\exp\int_0^T r_0(M)\dd t=\mathcal G(1/2,M_0)$.
Consequently the logarithm in \eqref{an02smcpv1:leadingratio}
is exactly $\int_0^T R(h)/h\dd t$.
Its lower bound follows by integrating on $[0,1]$ and its upper
bound by integrating the exponential upper bound on $[0,\infty)$.
Only the latter scalar integral extends to infinity; the
characteristic is not continued beyond $M=1/2$.
\end{proof}

The test characteristic need not carry mass, as in LOC's
same-field-reference formalism. This is an exact leading-system
example, not a claim that its atomic donor law equals OWPT's
initialized profile. The original-flow conclusion proved here is
Theorem~\ref{an02smcpv1:shortobstruction}; no unproved long-time
limit from the initialized population to this example is taken.

\section{Clock normalization, mean-mode audit, and stopping point}
Define, rather than suppress, the overlap mass coefficient
\begin{equation}\label{an02smcpv1:coefficient}
 A_q:=\frac{M_0}{e^{C_m}q^\sigma},\qquad 0<c\le A_q\le C.
\end{equation}
OWPT proves these bounds, not $A_q=1+o(1)$.
The exact handoff algebra is
\begin{equation}\label{an02smcpv1:handoff}
 \epsilon_mM_0^d=A_q^d q^{s+\sigma d}
 =A_q^d q^{10d-s}.
\end{equation}
The explicit exponentials in $C_m$ cancel. Equality without
$A_q^d$ applies to the mass scale $e^{C_m}q^\sigma$, not to the
actual mass as presently normalized. This preserves the proposed
$q$-power but retains an order-one shape prefactor.
For the defined leading logistic law only, the half-mass duration is
\begin{equation}\label{an02smcpv1:duration}
 T=\log\frac{1-M_0}{M_0}
 =\sigma\log(1/q)-C_m-\log A_q+\log(1-M_0).
\end{equation}
Adding $t_m$ gives $10\log(1/q)-\log A_q+O(q^\sigma)$.
This is algebra in the leading law, not a proved original-flow
duration or an assertion of its integrated mass defect.

LOC Lemma 3.1 and P's resident lemma use a \emph{stationary}
coherent leading center $(a,a)$ for every $M$ when $N=0$;
$J_M\preceq J_1$ concerns its collective attraction.
They do not specify a leading equilibrium moving with $M$ and
equal to $a_q$ at zero mass. The actual reference can move under
population forcing; a finite-$q$ coherent equilibrium, if used,
requires a separate definition and estimate. Nothing here resets
the actual reference or transfers shape through an additive error.

The stopped result is the central obstruction
\eqref{an02smcpv1:failedintegral}--\eqref{an02smcpv1:requiredfalse}.
The outer width $E_m$ is small if $C_m$ is large, but is not
$o_q(1)$ when $C_m$ is fixed. A shrinking-core restriction or
a new moving-overlap choice with $C_m\to\infty$, or a nonlinear
profile multiplier retaining the finite correction, would change
the requested contract. None is assumed or developed here.
In particular this note does not rule out an initialized theorem
with a bounded nonlinear shape factor, or prove that an endpoint-only
cancellation is impossible in that initialized flow. Such a
cancellation is not supplied by LOC's linear rate or same-field
ordering. Work stops without a half-mass original-flow theorem,
later radial-weight or off-chart propagation, or any later S1/S2
passage.

\begin{thebibliography}{9}\small
\bibitem{OWPT} OWPT, \emph{AN02 overlap window profile transfer v1}.
Labels \texttt{an02owptv1:jacobian}, \texttt{tangents},
\texttt{transfer}, \texttt{tail}, with that prefix throughout.
\bibitem{PROFILE} PROFILE, \emph{AN02 strong entry profile and tail
budget v1}. Prefix \texttt{inc:AN02:profile:v1:};
\texttt{root}, \texttt{entry}, \texttt{rightcircle}, \texttt{cost}.
\bibitem{LOC} LOC, \emph{AN03 local rate mean tracking v1}.
Prefix \texttt{inc:AN03:local:v1:}; \texttt{field},
\texttt{localineq}, \texttt{meanineq} (Lemma 3.1).
\bibitem{P} P, \emph{THEOREM A ANALYTICAL PROGRESS}, checkpoint 1.0.
\texttt{pa:leading}, \texttt{pa:masslaws}, \texttt{pa:resident},
\texttt{pa:spectra}, \texttt{pa:order}, and the exact early comparison.
\bibitem{ROADMAP} \emph{Unified analytic theory roadmap}, v1.3,
S0 assembly requirements and S1 step 2. Read-only source;
the roadmap's proposed route is not a proof premise.
\end{thebibliography}
\end{document}
